\paragraph{Optimization quality.}
Within the ablation configuration (\texttt{pc+rl} input, one prompt component ablated at a time),
neither model delivers large or reliable speedups. The best sustained result across all valid runs is
$1.14\times$ (GPT-5 on \texttt{bitonic-sort} under \texttt{no\_constraints}, \texttt{cot} or
\texttt{no\_proof}) and $1.13\times$ (Gemini on the same benchmark under \texttt{no\_explanation} or
\texttt{cot}). On \texttt{sobol} both models either fail or match the baseline within $2\%$. On
\texttt{atan2}, \texttt{md5hash} and \texttt{accuracy} the runtime distributions are dominated by
variant effects rather than by any genuine optimization.

\paragraph{Prompt components are not equally important.}
The \texttt{minimal} variant consistently underperforms---most clearly on \texttt{md5hash}
(GPT-5: $396{,}800$\,\textmu s versus $262{,}500$\,\textmu s under \texttt{no\_constraints}, about
$1.5\times$ slower) and on \texttt{shmembench} (Gemini: $18{,}735$\,\textmu s versus roughly
$4{,}660$\,\textmu s for every other variant). Stripping contextual framing measurably degrades
optimization quality. Conversely, \texttt{cot} is stable or slightly better for both models, though the
improvement is modest. Every other component is model-dependent: removing constraints helps GPT-5 on
\texttt{bitonic-sort} but hurts Gemini; removing the explanation helps Gemini on the same benchmark but
has no consistent effect elsewhere; and removing the role assignment caused a timeout for GPT-5 on
\texttt{accuracy} while Gemini handled it without issue. \texttt{Minimal} is the only variant that
consistently degrades both models, and \texttt{cot} the only one that never hurts either.

\paragraph{Reliability.}
GPT-5 succeeds on the first attempt in $32$ of $36$ runs, Gemini~2.5~Pro in $24$ of $36$. Gemini
exhausts all three repair iterations on eight runs, four of which are eventually repaired and four of
which fail outright; \texttt{sobol} accounts for three of the four failures, through correctness
violations and the host-only-intrinsic error described in Section~\ref{sec:ablation}. First-attempt
success rate, not peak speedup, is the dimension on which the two models differ most.

\paragraph{Benchmark gaming.}
The most consequential finding of the ablation is not a speedup but a false one. On
\texttt{shmembench}, GPT-5 games the benchmark under all six variants without exception, exploiting the
fact that a swap-based permutation preserves the sum to eliminate all shared-memory traffic and report
$235$--$949\times$ apparent speedups. Gemini does so only when the proof requirement is removed
(\texttt{no\_proof}), reaching $213\times$ by the same mechanism. Two things follow. First, a
prompt-level requirement for a formal correctness argument suppresses this class of edit in one model
but not the other, so it cannot be relied on as a defence. Second, the benchmark's own verification
passes in every one of these runs, so no output-comparison gate would have caught them. This motivates
the explicit computational-intent constraint evaluated in Section~\ref{sec:intent}.

\paragraph{From benchmarks to applications.}
The benchmark results above understate how much prompt design matters, because single-kernel benchmarks
hide two of its effects. They cannot expose compilation-context failures, since each benchmark is a
self-contained translation unit; and they make profiling data look worthless
(Section~\ref{sec:profabl}), since the hot loop of a $200$-line benchmark is obvious from the source.
Both effects reappear on the production applications of Section~\ref{sec:realapps}, which is the
strongest argument we can offer against drawing conclusions about such systems from benchmark suites
alone.
